\subsection{Frames and principal fibrations}

12.4.1. Let \(j \in \mathbf{J}^k(\mathbf{X}, \mathbf{Y})\) and let \(x = s(j)\), \(y = b(j)\). We say that \(j\) is invertible if there exists \(j' \in \mathbf{J}_y^k(\mathbf{Y}, \mathbf{X})_x\) such that \(j' \circ j = j_x^k(\mathrm{Id}_{\mathbf{X}})\) and \(j \circ j' = j_y^k(\mathrm{Id}_{\mathbf{Y}})\); the jet \(j'\) is then uniquely determined and is denoted \(j^{-1}\). If \(k = 0\), every jet is invertible. If \(k \geqslant 1\), the following conditions are equivalent:

\begin{enumerate}
\def\labelenumi{\alph{enumi})}
\item
  \(j\) is invertible;
\item
  the map \(\mathrm{T}(j)\colon\mathrm{T}_{x}(\mathrm{X})\rightarrow\mathrm{T}_{y}(\mathrm{Y})\) (cf. 12.3.4) is an isomorphism;
\item
  there exists an isomorphism \(g\) of class \(\mathbf{C}^r\) from an open neighborhood of \(x\) onto an open neighborhood of \(y\) such that \(j_x^k(g) = j\).
\end{enumerate}

12.4.2. Let E be a Banach space. We denote by \(\mathbf{GL}^{k}(\mathbf{E})\) the set of jets of order k from E to E that are invertible and have 0 as source and target; this is an open subset of \(\mathrm{J}_{0}^{k}(\mathrm{E},\mathrm{E})_{0}\). Equipped with the composition law of jets and with the manifold structure induced by that of the Banach space \(\mathrm{J}_{0}^{k}(\mathrm{E},\mathrm{E})_{0}=\mathrm{Q}^{k}(\mathrm{E},\mathrm{E})\), it is a group manifold of class \(C^{\omega}\).

One has \(\mathbf{GL}^0 (\mathbf{E}) = \{e\}\). The group \(\mathbf{GL}^1 (\mathbf{E})\) identifies, by means of T, with the group \(\mathbf{GL}(\mathbf{E})\) of automorphisms of E.

If \(k' \leqslant k\), the map \(r^{k, k'}: \mathbf{GL}^k(\mathbf{E}) \to \mathbf{GL}^{k'}(\mathbf{E})\) is a homomorphism of class \(\mathbf{C}^\omega\), and it is a surjective submersion. The map \(f \mapsto \mathrm{Id}_{\mathbf{E}} + f\) is an isomorphism of group manifolds from \(\mathrm{P}_k(\mathbf{E}, \mathbf{E})\) onto the kernel of \(r^{k, k-1}\).

12.4.3. Let E be a Banach space. For every \(x \in X\) , one calls an E-frame of order \(k\) of X at \(x\) every invertible element of \(J_0^k(E, X)_x\) . The set of E-frames of order \(k\) of X is an open submanifold \(R^k(E, X)\) of \(J_0^k(E, X)\) ; the map \(b\) has as restriction a morphism, still denoted \(b\) , of \(R^k(E, X)\) into X; likewise, if \(k' \leqslant k\) , we still denote by \(r^{k, k'}\) the morphism from \(R^k(E, X)\) to \(R^{k'}(E, X)\) , restriction of the map \(r^{k, k'}\) of 12.1.5.

12.4.4. Suppose that X is pure of type E (5.1.7). The group \(\mathbf{GL}^{k}(\mathbf{E})\) acts on the right on \(\mathrm{R}^{k}(\mathrm{E},\mathrm{X})\) by the law \((\rho,u)\mapsto\rho\circ u\), and the quadruple \(\lambda_{\mathbf{X}}=(\mathrm{R}^{k}(\mathbf{E},\mathrm{X}),\mathbf{GL}^{k}(\mathbf{E}),\mathbf{X},\mathbf{b})\) is a principal fibration (6.2.1) with structural group \(\mathbf{GL}^{k}(\mathbf{E})\), base X, and class \(C^{r-k}\).

If \(k' \leqslant k\), let H be the kernel of \(r^{k, k'} \colon \mathbf{GL}^{k}(\mathrm{E}) \to \mathbf{GL}^{k'}(\mathrm{E})\); the quadruple \((\mathbf{R}^k(\mathrm{E}, \mathrm{X}), \mathbf{H}, \mathbf{R}^{k'}(\mathrm{E}, \mathrm{X}), r^{k, k'})\) is a principal fibration of class \(\mathbf{C}^{r-k}\).

The manifold \(J^{k}(X,Y)\) is equipped with a structure of fiber bundle associated with \(\lambda_{X}\) (6.5.1): the typical fiber is \(J_{0}^{k}(E,Y)\), on which \(GL^{k}(E)\) acts on the left by the law \((u,j)\mapsto j\circ u^{-1}\); the frame map \(R^{k}(E,X)\times J_{0}^{k}(E,Y)\to J^{k}(X,Y)\) takes \((\rho,j)\) to \(j\circ\rho^{-1}\). The projection \(J^{k}(X,Y)\to X\) corresponding to this fiber-bundle structure is s.

12.4.5. Let F be a Banach space, and suppose Y is pure of type F. Then \(\mathrm{J}^{k}(\mathrm{X},\mathrm{Y})\) is equipped with a structure of fiber bundle associated with \(\lambda_{Y}\): the typical fiber is \(\mathrm{J}^{k}(\mathrm{X},\mathrm{F})_{0}\), on which \(\mathrm{GL}^{k}(\mathrm{F})\) acts on the left by the law \((v,j)\mapsto v\circ j\); the frame map is \((\sigma,j)\mapsto\sigma\circ j\); the projection \(\mathrm{J}^{k}(\mathrm{X},\mathrm{Y})\to\mathrm{Y}\) is b.

12.4.6. The hypotheses being those of 12.4.4 and 12.4.5, let \(\mu\) be the principal fibration

\[
(\mathbf {R} ^ {k} (\mathrm{E}, \mathrm{X}) \times \mathbf {R} ^ {k} (\mathrm{F}, \mathrm{Y}), \mathbf {G L} ^ {k} (\mathrm{E}) \times \mathbf {G L} ^ {k} (\mathrm{F}), \mathrm{X} \times \mathrm{Y}, \boldsymbol {b} \times \boldsymbol {b}),
\]

the product of \(\lambda_{\mathbf{X}}\) and \(\lambda_{\mathbf{Y}}\). Then \(\mathrm{J}^{k}(\mathrm{X},\mathrm{Y})\) is equipped with a structure of fiber bundle associated with \(\mu\): the typical fiber is \(\mathrm{J}_{0}^{k}(\mathrm{E},\mathrm{F})_{0}\), on which \(\mathbf{GL}^{k}(\mathrm{E})\times \mathbf{GL}^{k}(\mathrm{F})\) acts on the left by the law \(((u,v),j)\mapsto v\circ j\circ u^{-1}\); the frame map is \(((\rho ,\sigma),j)\mapsto \sigma \circ j\circ \rho^{-1}\); the projection \(\mathrm{J}^k (\mathrm{X},\mathrm{Y})\to \mathrm{X}\times \mathrm{Y}\) is \((s,b)\).

